\documentclass[12pt,a4paper]{report}

\usepackage[margin=1in]{geometry}
\usepackage{setspace}
\usepackage{titlesec}
\usepackage{graphicx}
\usepackage{booktabs}
\usepackage{longtable}
\usepackage{array}
\usepackage{enumitem}
\usepackage{hyperref}
\usepackage{amsmath,amssymb}
\usepackage{tocloft}

\setstretch{1.2}
\hypersetup{colorlinks=true,linkcolor=black,urlcolor=blue,citecolor=black}
\setcounter{secnumdepth}{3}
\setcounter{tocdepth}{2}

\begin{document}

% -----------------------------------------------------------------
% PART A: COVER PAGE
% -----------------------------------------------------------------
\begin{titlepage}
\centering
{\Large INDUSTRY INTERNSHIP REPORT / THESIS\par}
\vspace{0.8cm}
{\normalsize Submitted in partial fulfillment of requirement for the award of degree of\par}
\vspace{0.5cm}
{\Large Bachelor of Technology in Computer Science and Engineering\par}
\vspace{0.8cm}
{\large by\par}
\vspace{0.3cm}
{\Large Punyaa Dixit\par}
\vspace{0.8cm}
{\large Industry / Organization Guide\par}
{\large Dr. / Mr. / Mrs. \rule{7cm}{0.4pt}\par}
\vspace{0.6cm}
{\large at\par}
{\Large Local Large Language Model Hallucination Detection Project\par}
\vspace{0.6cm}
{\large Institute Guide (from College)\par}
{\large Dr. Apeksha Sakhare\par}
{\large Assistant / Associate Professor, Department of CSE\par}
\vspace{0.8cm}
{\large May 2026\par}
\vfill
{\large Department of Computer Science and Engineering\par}
{\large G H Raisoni College of Engineering, Nagpur\par}
{\normalsize An Empowered Autonomous Institute affiliated to RTMNU, Nagpur\par}
\end{titlepage}

% -----------------------------------------------------------------
% PART A: INNER PAGE
% -----------------------------------------------------------------
\begin{titlepage}
\centering
{\Large INDUSTRY INTERNSHIP REPORT / THESIS\par}
\vspace{0.8cm}
{\normalsize Submitted in partial fulfillment of requirement for the award of degree of\par}
\vspace{0.5cm}
{\Large Bachelor of Technology in Computer Science and Engineering\par}
\vspace{0.8cm}
{\large by\par}
\vspace{0.3cm}
{\Large Punyaa Dixit\par}
\vspace{0.8cm}
{\large Industry / Organization Guide\par}
{\large Dr. / Mr. / Mrs. \rule{7cm}{0.4pt}\par}
\vspace{0.6cm}
{\large at\par}
{\Large Local Large Language Model Hallucination Detection Project\par}
\vspace{0.6cm}
{\large Institute Guide (from College)\par}
{\large Dr. Apeksha Sakhare\par}
{\large Assistant / Associate Professor, Department of CSE\par}
\vspace{0.8cm}
{\large May 2026\par}
\vfill
{\large Department of Computer Science and Engineering\par}
{\large G H Raisoni College of Engineering, Nagpur\par}
\vspace{0.3cm}
{\small Copyright \textcopyright{} G H Raisoni College of Engineering, Nagpur, 2026\par}
\end{titlepage}

% -----------------------------------------------------------------
% PART A: DECLARATION
% -----------------------------------------------------------------
\chapter*{Declaration}
\addcontentsline{toc}{chapter}{Declaration}
I hereby declare that this Industry Internship Report / Thesis entitled
\textbf{"A Multi-Stage Framework for Detecting Factual Hallucinations in Large Language Model Outputs"}
submitted in partial fulfillment of the requirements for the award of
Bachelor of Technology in Computer Science and Engineering has been carried out by me.

The work presented in this report is original and has not been submitted, in whole or in part,
for the award of any other degree or diploma at this or any other institution or university.

I also assign to G H Raisoni College of Engineering, Nagpur all rights under copyright that may
exist in and to the above work and any revised or expanded derivative works based on this work.

\vspace{1cm}
\begin{tabular}{p{0.45\textwidth} p{0.45\textwidth}}
Name of Student: Punyaa Dixit & Signature: \rule{5cm}{0.4pt} \\
Mobile No.: \rule{5cm}{0.4pt} & Date: \rule{5cm}{0.4pt} \\
Email ID: \rule{5cm}{0.4pt} & Place: Nagpur \\
\end{tabular}

% -----------------------------------------------------------------
% PART A: CERTIFICATE
% -----------------------------------------------------------------
\chapter*{Certificate}
\addcontentsline{toc}{chapter}{Certificate}
This is to certify that the Industry Internship Report / Thesis entitled
\textbf{"A Multi-Stage Framework for Detecting Factual Hallucinations in Large Language Model Outputs"}
has been carried out under our supervision and guidance by
\textbf{Punyaa Dixit}
for the award of Degree of Bachelor of Technology in Computer Science and Engineering.

The work submitted is comprehensive, complete, and fit for evaluation.

\vspace{1.2cm}
\begin{tabular}{p{0.45\textwidth} p{0.45\textwidth}}
Dr. / Mr. / Mrs. \rule{5cm}{0.4pt} & Dr. Apeksha Sakhare \\
Industry / Organization Guide & Institute Guide \\
Designation: \rule{5cm}{0.4pt} & Assistant / Associate Professor, CSE \\
Organization: \rule{5cm}{0.4pt} & GHRCE, Nagpur \\
\end{tabular}

\vspace{1.0cm}
\noindent\textbf{Note:} Attach the official industry certificate(s) on organization letterhead with
reference/document number, internship duration, project details, and stipend/sponsorship details,
as required by the departmental guideline.

% -----------------------------------------------------------------
% PART A: ACKNOWLEDGEMENT
% -----------------------------------------------------------------
\chapter*{Acknowledgement}
\addcontentsline{toc}{chapter}{Acknowledgement}
I express my sincere gratitude to my institute guide, \textbf{Dr. Apeksha Sakhare},
for continuous guidance, constructive feedback, and support throughout this work.
I also thank the Department of Computer Science and Engineering, G H Raisoni College of Engineering,
Nagpur, for providing an encouraging academic environment.

I acknowledge the open-source community and research ecosystem, including Hugging Face,
spaCy, Sentence-Transformers, and Wikipedia contributors, whose tools and knowledge sources
were instrumental in developing and evaluating this hallucination detection framework.

% -----------------------------------------------------------------
% PART B: ROMAN NUMBERED PRELIMS
% -----------------------------------------------------------------
\pagenumbering{roman}

\chapter*{Abstract}
\addcontentsline{toc}{chapter}{Abstract}
Large Language Models (LLMs) are widely adopted in content generation and decision-support workflows,
but they often produce hallucinations: fluent statements that are factually incorrect or unverifiable.
This thesis presents a multi-stage framework for detecting factual hallucinations in LLM outputs.
The framework includes claim extraction, evidence retrieval, verification using natural language
inference, and a risk-based hallucination scoring layer.

A modular Python implementation was developed using spaCy, transformer-based inference models,
and Wikipedia-backed retrieval. Evaluation on curated scientific, historical, geographic,
and intentionally hallucinated samples demonstrated strong performance, with
\textbf{82.1\% overall claim-level accuracy} and \textbf{87.3\% precision for contradicted claims}.
The system is designed for practical deployment and includes extensible interfaces for local
and service-based model backends. This work contributes a practical, explainable verification
pipeline that can improve trust and safety in real-world LLM applications.

\tableofcontents
\listoffigures
\listoftables

\chapter*{List of Symbols}
\addcontentsline{toc}{chapter}{List of Symbols}
\begin{longtable}{p{0.2\textwidth} p{0.72\textwidth}}
$T$ & Input LLM-generated text \\
$C$ & Set of extracted factual claims \\
$E$ & Evidence set retrieved for a claim \\
$V(c)$ & Verification label for claim $c$ \\
$H$ & Overall hallucination score for a document \\
$\theta$ & Support threshold in verification aggregation \\
\end{longtable}

\chapter*{List of Publications / Competitions / Patents}
\addcontentsline{toc}{chapter}{List of Publications / Competitions / Patents}
\begin{itemize}[leftmargin=1.5em]
\item Research paper: \textit{A Multi-Stage Framework for Detecting Factual Hallucinations in Large Language Model Outputs}.
\item Add published paper details / competition participation / patent filing details here, if applicable.
\end{itemize}

% -----------------------------------------------------------------
% PART C: MAIN BODY WITH ARABIC NUMBERS
% -----------------------------------------------------------------
\clearpage
\pagenumbering{arabic}

\chapter*{Index}
\addcontentsline{toc}{chapter}{Index}
The chapter sequence of this report follows the institutional industry report guideline.

\chapter{Introduction to Company / Organization Context}
This thesis was executed as an industry-aligned applied AI project in the domain of
\textbf{trustworthy large language model systems}. The project context is software-product
engineering, where generated content must be validated before deployment in user-facing workflows.

The practical problem addressed is: \textit{how to detect factual hallucinations in generated text
without retraining the base LLM}. The developed solution is a modular post-hoc verification pipeline
suitable for integration with local and cloud model ecosystems.

\section{About the Project Environment}
The implementation environment consisted of Python 3.12, model inference libraries,
NLP toolkits, and a lightweight web interface for interactive analysis.
The workflow included data preparation, pipeline experimentation,
error analysis, and reproducible reporting.

\section{Historical Background}
As LLM capabilities expanded rapidly between 2022 and 2026,
concerns about hallucinated content also increased. Prior studies showed that fluency does not imply truthfulness.
This motivated development of independent verification systems that can audit generated claims against trusted evidence.

\section{Location and Operational Structure}
The project was carried out within the academic and research setup of
G H Raisoni College of Engineering, Nagpur, with guidance from faculty mentorship.
Operationally, work followed iterative phases: literature review, architecture design,
implementation, evaluation, and final documentation.

\section{Vision and Mission Alignment}
\textbf{Vision:} Improve reliability and safety of AI-generated content through transparent verification.

\textbf{Mission:}
\begin{itemize}[leftmargin=1.5em]
\item Build a practical framework for detecting hallucinations in LLM outputs.
\item Provide interpretable, claim-level evidence for user trust.
\item Support deployment in educational, research, and industry assistant applications.
\end{itemize}

\section{Product / Solution Developed}
The resulting product is a configurable hallucination detector with these major modules:
claim extraction, evidence retrieval, NLI-based verification, and risk scoring.
A web interface and script-based APIs are provided for integration and testing.

\chapter{Case Study / Project Work}
\section{Introduction}
The project investigates factual hallucination detection as a structured verification problem.
Instead of modifying language model generation, the system evaluates completed model outputs
and quantifies factual risk.

\section{Problem Identification}
LLM outputs can include fabricated entities, wrong dates, incorrect relationships,
and unverifiable claims. In high-impact domains, such errors reduce trust and can
cause downstream operational risk.

\section{Objectives}
\begin{itemize}[leftmargin=1.5em]
\item Extract atomic factual claims from free-form generated text.
\item Retrieve high-relevance supporting evidence from trusted sources.
\item Classify each claim as supported, contradicted, or unverifiable.
\item Produce an overall risk score and interpretable report.
\item Validate effectiveness across mixed-domain evaluation data.
\end{itemize}

\section{Work Carried Out}
\subsection{Pipeline Design}
A four-stage architecture was designed:
\begin{enumerate}[leftmargin=1.5em]
\item Claim extraction using linguistic and entity-aware methods.
\item Evidence retrieval using query-focused lookup and semantic filtering.
\item Verification via transformer-based NLI inference.
\item Aggregated hallucination scoring and risk labeling.
\end{enumerate}

\subsection{Hardware and Software Setup}
\begin{itemize}[leftmargin=1.5em]
\item Hardware: Standard development workstation; optional GPU acceleration.
\item Software: Python 3.12, Transformers, Sentence-Transformers, spaCy, NLTK, Wikipedia API.
\item Interface: Streamlit-based web front end for interactive claim analysis.
\end{itemize}

\subsection{Solution Implementation}
The detector was implemented as modular classes in Python for maintainability:
\textit{ClaimExtractor}, \textit{EvidenceRetriever}, \textit{InferenceModel}, and \textit{HallucinationScorer}.
This design supports testing and replacement of individual components.

\section{Solution Provided}
The final solution is an end-to-end hallucination assessment engine that accepts generated text and returns:
\begin{itemize}[leftmargin=1.5em]
\item Extracted factual claims.
\item Matched evidence snippets and relevance confidence.
\item Claim-level verification labels (Supported/Contradicted/Unverifiable).
\item Document-level hallucination score and risk band.
\end{itemize}

\section{Calculation Done}
Claim-level risk scores are mapped as:
\begin{equation}
conf(v_i)=
\begin{cases}
0.0, & v_i=\text{SUPPORTED} \\
0.5, & v_i=\text{UNVERIFIABLE} \\
1.0, & v_i=\text{CONTRADICTED}
\end{cases}
\end{equation}

Overall hallucination score:
\begin{equation}
H=\frac{1}{n}\sum_{i=1}^{n} conf(v_i)
\end{equation}
where $n$ is the number of extracted claims.

Relevance scoring between claim and evidence embedding vectors:
\begin{equation}
relevance(c,e)=\frac{emb(c)\cdot emb(e)}{\|emb(c)\|\|emb(e)\|}
\end{equation}

\section{Results and Conclusion}
Evaluation outcomes from the final selected research paper are summarized below.

\begin{table}[h!]
\centering
\caption{Verification Performance Summary}
\begin{tabular}{lccc}
\toprule
Status & Precision & Recall & F1 \\
\midrule
Supported & 81.4\% & 78.6\% & 80.0\% \\
Contradicted & 87.3\% & 83.2\% & 85.2\% \\
Unverifiable & 74.8\% & 79.3\% & 77.0\% \\
\midrule
Overall & 82.1\% & 80.4\% & 81.2\% \\
\bottomrule
\end{tabular}
\end{table}

\begin{table}[h!]
\centering
\caption{Accuracy by Content Domain}
\begin{tabular}{lcc}
\toprule
Domain & Claim Accuracy & Document Accuracy \\
\midrule
Scientific Facts & 86.3\% & 92.0\% \\
Historical Events & 79.2\% & 86.0\% \\
Geographic Information & 82.7\% & 88.0\% \\
Hallucinated Content & 81.4\% & 88.0\% \\
\bottomrule
\end{tabular}
\end{table}

The framework achieved reliable detection of contradicted claims and practical end-to-end throughput,
showing suitability for real-world AI verification workflows.

\chapter{References}
\begin{enumerate}[leftmargin=1.5em]
\item L. Huang, W. Yu, W. Ma et al., "A Survey on Hallucination in Large Language Models: Principles, Taxonomy, Challenges, and Open Questions," ACM Transactions on Information Systems, vol. 43, no. 2, 2025.
\item M. Cleti and P. Jano, "Hallucinations in LLMs: Types, Causes, and Approaches for Enhanced Reliability," 2024.
\item S. Farquhar, J. Kossen, L. Kuhn, and Y. Gal, "Detecting hallucinations in large language models using semantic entropy," Nature, vol. 630, pp. 625-630, 2024.
\item A. T. Kalai, O. Nachum, S. S. Vempala, and E. Zhang, "Why Language Models Hallucinate," 2025.
\item J. Thorne, A. Vlachos, C. Christodoulopoulos, and A. Mittal, "FEVER: a Large-scale Dataset for Fact Extraction and VERification," Proceedings of NAACL-HLT, pp. 809-819, 2018.
\item Y. Liu, M. Ott, N. Goyal et al., "RoBERTa: A Robustly Optimized BERT Pretraining Approach," arXiv preprint arXiv:1907.11692, 2019.
\item P. Lewis, E. Perez, A. Piktus et al., "Retrieval-Augmented Generation for Knowledge-Intensive NLP Tasks," Advances in Neural Information Processing Systems, vol. 33, pp. 9459-9474, 2020.
\item X. Wang, J. Wei, D. Schuurmans et al., "Self-Consistency Improves Chain of Thought Reasoning in Language Models," International Conference on Learning Representations, 2023.
\item N. Reimers and I. Gurevych, "Sentence-BERT: Sentence Embeddings using Siamese BERT-Networks," Proceedings of EMNLP-IJCNLP, pp. 3982-3992, 2019.
\item M. Honnibal, I. Montani, S. Van Landeghem, and A. Boyd, "spaCy: Industrial-Strength Natural Language Processing in Python," 2020. [Online]. Available: \url{https://spacy.io}
\end{enumerate}

\chapter{Appendices}
\section*{Appendix A: Image/Photo List of Project Work}
\addcontentsline{toc}{section}{Appendix A: Image/Photo List of Project Work}
\begin{itemize}[leftmargin=1.5em]
\item Add screenshots of system architecture, web interface, and result dashboards.
\item Add group photo with guide / co-guide (if applicable).
\end{itemize}

\section*{Appendix B: Industry/Organization Engagement Evidence}
\addcontentsline{toc}{section}{Appendix B: Industry/Organization Engagement Evidence}
\begin{itemize}[leftmargin=1.5em]
\item Add photos from organization visits / interactions with supervisors.
\item Attach signed industry certificates with reference/document number.
\item Add cost-savings or implementation impact certificate, if available.
\end{itemize}

\section*{Appendix C: Publications / Competitions / Innovation Artifacts}
\addcontentsline{toc}{section}{Appendix C: Publications / Competitions / Innovation Artifacts}
\begin{itemize}[leftmargin=1.5em]
\item Add publication acceptance/email proof (if available).
\item Add project competition participation/winning certificates.
\item Add patent filing copy (if filed).
\end{itemize}

\end{document}
